\subsection{分次代数的 Poincaré 级数}

设 K 是一个有单位元且不退化为 0 的交换环。设 M 是一个 Z 型分次 K-模，\(M_{n}\) 是 M 中次数为 n 的齐次元素所成的集合。假设每个 \(M_{n}\) 都是自由的且有限型的。于是秩 \(\operatorname{rk}_{\mathbf{K}}(\mathbf{M}_{n})\) 对所有 n 有定义（《交换代数》第 II 章，§5，第 3 号）。

定义 1. 若存在 \(n_0 \in \mathbf{Z}\) 使 \(M_n = 0\) （对于 \(n \leqslant n_0\) ），则形式级数 \(\sum_{n \geqslant n_0} \mathrm{rk}_K(M_n) T^n\) （\(\mathbf{Q}((T))\) 的一个元素）称为 \(M\) 的 Poincaré 级数，记为 \(P_M(T)\) 。

设 \(M'\) 是另一个 Z 型分次 K-模，\((M_n')_{n \in \mathbf{Z}}\) 是它的分次。假设 \(M_n'\) 对小于某个数的 n 为零。则

\[
\mathrm{P} _ {\mathrm{M} \oplus \mathrm{M} ^ {\prime}} (\mathrm{T}) = \mathrm{P} _ {\mathrm{M}} (\mathrm{T}) + \mathrm{P} _ {\mathrm{M} ^ {\prime}} (\mathrm{T})\tag{1}
\]

而若给 \(\mathbf{M} \otimes_{\mathbf{K}} \mathbf{M}'\) 赋予总分次（《代数》第 II 章，§11，第 5 号），则

\[
\mathrm{P} _ {\mathrm{M} \otimes \mathrm{M} ^ {\prime}} (\mathrm{T}) = \mathrm{P} _ {\mathrm{M}} (\mathrm{T}) \mathrm{P} _ {\mathrm{M} ^ {\prime}} (\mathrm{T}).\tag{2}
\]

命题 1. 设 \(S = \bigoplus_{n \geqslant 0} S_n\) 是一个交换分次 K-代数，具有一个由齐次且代数无关的元素组成的生成元组 \((x_1, x_2, \ldots, x_m)\) 。设 \(d_i\) 是 \(x_i\) 的次数，并假设 \(d_i > 0\) （对所有 \(i\) ）。则诸 \(S_n\) 在 \(K\) 上是自由的且有限秩的，并且

\[
\mathrm{P} _ {\mathrm{S}} (\mathrm{T}) = \prod_ {i = 1} ^ {m} (1 - \mathrm{T} ^ {d _ {i}}) ^ {- 1}.\tag{3}
\]

事实上，S 可与张量积 \(K[x_{1}] \otimes \cdots \otimes K[x_{m}]\) 等同，后者赋予总分次。\(K[x_{i}]\) 的 Poincaré 级数是

\[
\sum_ {n \geqslant 0} \mathrm{T} ^ {n d _ {i}} = (1 - \mathrm{T} ^ {d _ {i}}) ^ {- 1},
\]

于是只需应用 (2)。

在命题 1 的假设下，我们将称 S 是 K 上的一个分次多项式代数。

推论. 诸次数 \(d_{i}\) 由 S 确定，至多相差次序。

事实上，\(\mathrm{P}_{\mathrm{S}}(\mathrm{T})\) 的逆是多项式 \(\mathrm{N}(\mathrm{T}) = \prod_{i=1}(1 - \mathrm{T}^{d_i})\) ，从而是唯一确定的。若 q 是一个整数 \(\geqslant 1\) 且 \(\zeta \in C\) 是一个本原 q 次单位根，则根 \(\zeta\) 在 \(\mathrm{N}(\mathrm{T})\) 中的重数等于 \(d_{i}\) 中是 q 的倍数的个数。当 q 充分大时这个数为零。于是等于 q 的 \(d_{i}\) 的个数由递降归纳唯一确定。

诸整数 \(d_{i}\) 称为 S 的特征次数。当 K 是域时，它们的个数等于 S 在 K 上的超越次数；在一般情形我们也把它称为 S 在 K 上的超越次数。它是多项式 N(T) 的根 1 的重数。

设 \(S = \bigoplus_{n \geqslant 0} S_n\) 是一个交换分次 K-代数，\(R = \bigoplus_{n \geqslant 0} R_n\) 是 S 的一个分次子代数。假设每个 \(R_n\) 都是自由的且有限型的，且 R-模 S 具有一个由齐次元素 \(z_1, z_2, \ldots, z_N\) （次数分别为 \(f_1, f_2, \ldots, f_N\) ）组成的有限基。于是，若用 M 表示分次 K-模 \(\sum_{j=1}^{N} Kz_j\) ，则分次 K-模 S 同构于 \(R \otimes_K M\) ，从而每个 \(S_n\) 都是自由的且有限型的，并且

\[
P _ {S} (T) = P _ {M} (T) P _ {R} (T) = \left(\sum_ {j = 1} ^ {N} T ^ {f _ {j}}\right) P _ {R} (T).\tag{4}
\]

命题 2. 沿用前面的记号，并假设 \(S\) 与 \(R\) 是分次多项式代数。

\begin{enumerate}
\def\labelenumi{(\roman{enumi})}
\item
  R 与 S 有相同的超越次数 \(r\) （在 \(K\) 上）。
\item
  设 \(p_1, \ldots, p_r\) （分别为 \(q_1, \ldots, q_r\) ）是 \(S\) （分别为 \(R\) ）的特征次数。则
\end{enumerate}

\[
\prod_ {i = 1} ^ {r} \left(1 - \mathrm{T} ^ {q _ {i}}\right) = \left(\sum_ {j = 1} ^ {\mathrm{N}} \mathrm{T} ^ {f _ {j}}\right) \prod_ {i = 1} ^ {r} \left(1 - \mathrm{T} ^ {p _ {i}}\right).
\]

\begin{enumerate}
\def\labelenumi{(\roman{enumi})}
\setcounter{enumi}{2}
\tightlist
\item
  \(Np_{1}p_{2}\ldots p_{r}=q_{1}q_{2}\ldots q_{r}.\)
\end{enumerate}

公式 (4) 首先表明根 1 的重数在多项式 \(\mathrm{P}_{\mathrm{S}}(\mathrm{T})^{-1}\) 与 \(\mathrm{P}_{\mathrm{R}}(\mathrm{T})^{-1}\) 中相同，再考虑到 (3)，就同时证明了 (i) 与 (ii)。

由 (ii) 得到

\[
\prod_ {i = 1} ^ {r} \left(1 + \mathrm{T} + \mathrm{T} ^ {2} + \dots + \mathrm{T} ^ {q _ {i} - 1}\right) = \left(\sum_ {j = 1} ^ {\mathrm{N}} \mathrm{T} ^ {f _ {j}}\right) \prod_ {i = 1} ^ {r} \left(1 + \mathrm{T} + \mathrm{T} ^ {2} + \dots + \mathrm{T} ^ {p _ {i} - 1}\right).
\]

在此关系中置 \(\mathrm{T} = 1\) 即得 (iii)。

注. 设 \(S = K[X_1, \ldots, X_n]\) 是 \(K\) 上的一个分次多项式代数，\(d_i\) 是 \(X_i\) 的次数，\(F(X_1, \ldots, X_n)\) 是一个次数为 \(m\) 的 \(S\) 的齐次元素。则

\[
\sum_ {i = 1} ^ {n} d _ {i} X _ {i} \frac {\partial F}{\partial X _ {i}} = m F.\tag{5}
\]

事实上，立即可见把每个次数为 p 的齐次元素 z 变为 pz 的、从 S 到 S 的 K-线性映射 D 是 S 的一个求导。于是

\[
m \mathrm{F} (\mathrm{X} _ {1}, \dots , \mathrm{X} _ {n}) = \mathrm{D} (\mathrm{F} (\mathrm{X} _ {1}, \dots , \mathrm{X} _ {n})) = \sum_ {i = 1} ^ {n} \mathrm{D} (\mathrm{X} _ {i}) \frac {\partial \mathrm{F}}{\partial \mathrm{X} _ {i}} = \sum_ {i = 1} ^ {n} d _ {i} \mathrm{X} _ {i} \frac {\partial \mathrm{F}}{\partial \mathrm{X} _ {i}}.
\]

